% Part: first-order-logic
% Chapter: sequent-calculus
% Section: provability-consistency

\documentclass[../../../include/open-logic-section]{subfiles}

\begin{document}

\iftag{FOL}
      {\olfileid[tr]{fol}{seq}{prv}}
      {\olfileid[tr]{pl}{seq}{prv}}

\olsection{\usetoken{S}{derivability} ve Tutarlılık}

Şimdi !!{derivability} bağıntısının bir dizi özelliğini ortaya koyacağız. Bu
özelliklerin her biri kendi başına ilginçtir ve tamlık teoreminin ispatında rol
oynayacaktır.

\begin{prop}\ollabel{prop:provability-contr}
  $\Gamma \Proves !A$ ve $\Gamma\cup\{!A\}$ tutarsızsa $\Gamma$ tutarsızdır.
\end{prop}

\begin{proof}
Sonlu $\Gamma_0,\Gamma_1\subseteq\Gamma$ kümeleri vardır ve $\Log{LK}$ içinde
$\Gamma_0\Sequent !A$ ile $!A,\Gamma_1\Sequent\quad$ sekantlarının ikisini de
!!{derive}s mümkündür. $\Gamma_0\Sequent !A$ sekantının $\Log{LK}$-!!{derivation}{}i $\pi_0$,
$!A,\Gamma_1\Sequent\quad$ sekantının $\Log{LK}$-!!{derivation}{}i ise $\pi_1$
olsun. O zaman aşağıdakini !!{derive} mümkündür:
\begin{prooftree}
\AxiomC{}
\RightLabel{$\pi_0$}
\Deduce$ \Gamma_0 \fCenter !A $
\AxiomC{}
\RightLabel{$\pi_1$}
\Deduce$!A, \Gamma_1 \fCenter $
\RightLabel{\Cut}
\BinaryInf$ \Gamma_0,\Gamma_1 \fCenter $
\end{prooftree}

$\Gamma_0\subseteq\Gamma$ ve $\Gamma_1\subseteq\Gamma$ olduğundan
$\Gamma_0\cup\Gamma_1\subseteq\Gamma$; dolayısıyla $\Gamma$ tutarsızdır.
\end{proof}

\begin{prop}
\ollabel{prop:prov-incons}
$\Gamma \Proves !A$ ancak ve ancak $\Gamma\cup\{\lnot !A\}$ tutarsızsa
geçerlidir.
\end{prop}

\begin{proof}
Önce $\Gamma\Proves !A$ olduğunu varsayalım; yani sonlu bir
$\Gamma_0\subseteq\Gamma$ için $\Gamma_0\Sequent !A$ sekantının
!!a{derivation}{}i $\pi_0$ vardır. Bir $\LeftR{\lnot}$ kuralı ekleyerek
$\lnot !A,\Gamma_0\Sequent\quad$ sekantının !!a{derivation}{}ini elde ederiz;
yani $\Gamma\cup\{\lnot !A\}$ tutarsızdır.

$\Gamma\cup\{\lnot !A\}$ tutarsızsa sonlu bir $\Gamma_1\subseteq\Gamma$
için $\lnot !A,\Gamma_1\Sequent\quad$ sekantının !!a{derivation}{}i $\pi_1$
vardır. Aşağıdaki, $\Gamma_1\Sequent !A$ sekantının !!a{derivation}{}idir:
\begin{prooftree}
  \Axiom$!A \fCenter !A$
  \RightLabel{\RightR{\lnot}}
  \UnaryInf$\fCenter !A, \lnot !A$
  \AxiomC{}
  \RightLabel{$\pi_1$}
  \Deduce$\lnot !A, \Gamma_1 \fCenter$
  \RightLabel{\Cut}
  \BinaryInf$\Gamma_1 \fCenter !A$
\end{prooftree}
\end{proof}

\begin{prob}
$\Gamma\Proves\lnot !A$ olmasının $\Gamma\cup\{!A\}$ kümesinin tutarsız
olmasına denk olduğunu ispatlayın.
\end{prob}

\begin{prop}\ollabel{prop:explicit-inc}
  $\Gamma\Proves !A$ ve $\lnot !A\in\Gamma$ ise $\Gamma$ tutarsızdır.
\end{prop}

\begin{proof}
  $\Gamma\Proves !A$ ve $\lnot !A\in\Gamma$ olduğunu varsayalım. Sonlu bir
  $\Gamma_0\subseteq\Gamma$ için $\Gamma_0\Sequent !A$ sekantının
  !!a{derivation}{}i $\pi$ vardır.
$\lnot !A,\Gamma_0\Sequent\quad$ sekantı da !!{derivable}{}dir:
  \begin{prooftree}
    \AxiomC{}
    \RightLabel{$\pi$}
    \Deduce$\Gamma_0 \fCenter !A$
    \Axiom$!A \fCenter !A$
    \RightLabel{\LeftR{\lnot}}
    \UnaryInf$\lnot !A, !A \fCenter$
    \RightLabel{\LeftR{\Exchange}}
    \UnaryInf$!A, \lnot !A \fCenter$
    \RightLabel{\Cut}
    \BinaryInf$\Gamma_0, \lnot !A \fCenter$
  \end{prooftree}
  $\lnot !A\in\Gamma$ ve $\Gamma_0\subseteq\Gamma$ olduğundan bu,
  $\Gamma$'nın tutarsız olduğunu gösterir.
\end{proof}

\begin{prop}\ollabel{prop:provability-exhaustive}
  $\Gamma\cup\{!A\}$ ile $\Gamma\cup\{\lnot !A\}$ kümelerinin ikisi de
  tutarsızsa $\Gamma$ tutarsızdır.
\end{prop}

\begin{proof}
Sonlu $\Gamma_0\subseteq\Gamma$ ve $\Gamma_1\subseteq\Gamma$ kümeleri ile
sırasıyla $!A,\Gamma_0\Sequent\quad$ ve $\lnot !A,\Gamma_1\Sequent\quad$
sekantlarının $\Log{LK}$-!!{derivation}s $\pi_0$ ve $\pi_1$ vardır. O zaman
aşağıdakini !!{derive} mümkündür:
\begin{prooftree}
\AxiomC{}
\RightLabel{$\pi_0$}
\Deduce$ !A, \Gamma_0 \fCenter $
\RightLabel{\RightR{\lnot}}
\UnaryInf$ \Gamma_0 \fCenter \lnot !A$
\AxiomC{}
\RightLabel{$\pi_1$}
\Deduce$\lnot !A, \Gamma_1 \fCenter  $
\RightLabel{\Cut}
\BinaryInf$ \Gamma_0, \Gamma_1 \fCenter $
\end{prooftree}
$\Gamma_0\subseteq\Gamma$ ve $\Gamma_1\subseteq\Gamma$ olduğundan
$\Gamma_0\cup\Gamma_1\subseteq\Gamma$. Dolayısıyla $\Gamma$ tutarsızdır.
\end{proof}

\end{document}
